\scanpage{372}
\mspage{\deleted{\unsure{110}}\unsure{111|110}}
% archive sheet number: 371
semble la façon la plus simple, et peut-être unique, de prouver que l'image de $C_\infty$ dans $\mathrm{Hot}$ ne dépend pas, à isom. unique près, des choix faits pour construire $C_\infty$. Sans compter que dans le cas \unsure{motivant}, $C=\underline{\mathrm{Ch}}_{\unsure{\Sigma|Z}}(X)$, ou $\underline{\mathrm{Ch}}_{\Sigma}(X;a,b)$, $T=\unsure{\underline{\mathrm{Tg}}|\underline{T_g}}$, \inserted{c'est} la catégorie \unsure{des|de} chemins $\deleted{\unsure{C=\mathrm{Ch}}}$ $C$ qui donne lieu à une \deleted{\unsure{seule}} loi de composition \inserted{des chemins} impeccable, fonctorielle, associative et tout, alors qu'il n'en est rien pour $C_\infty$. On ne trouve pas \deleted{\unsure{de}} une loi de composition
\[
C_\infty\mathbin{\unsure{\times|<}}C_\infty\mathrel{\unsure{-|\to}}C_\infty
\]
dans \deleted{\unsure{Hot,}} $\mathrm{Cat}$, \deleted{\unsure{ou}} ni même dans $\mathrm{Cat\,hot}$ ou dans $\mathrm{Hot}$, ou plutôt, si on arrive \inserted{\unsure{pourtant|péniblement}} à \deleted{définir} construire \inserted{une flèche} $C_\infty\times C_\infty\longrightarrow C_\infty$ dans $\mathrm{Cat}$, c'est encore \inserted{laborieusement,} à \unsure{coup|coups} de choix $\pm$ arbitraires, \unsure{cran par cran|cas par cas}, et sans savoir si ce qu'on obtient a une \deleted{\unsure{sig}} propriété d'\unsure{unicité}, ne serait-ce que dans $\mathrm{Hot}$ ; on ne trouve une interprétation géométrique \unsure{tangible|intelligible} \inserted{de cette composition} que \deleted{\unsure{par le}} dans $\mathrm{Hot}_0^{\unsure{\wedge|1}}$ (où on a une véritable loi de groupoïde), ou dans $H_T$ en tant qu'objet de $\mathrm{Ind\,Cat\,hot}$
